\documentclass{article}
\usepackage{PRIMEarxiv} % Core package for the PrimeAI Template
\usepackage{amsmath, amssymb, amsthm, bm, dcolumn} % Add amsthm here for the proof environment
\usepackage[numbers,sort&compress]{natbib} % Natbib for citations
\usepackage{graphicx} % For high-quality images
\usepackage{multicol} % Multi-column figures/tables
\usepackage{hyperref} % Hyperlinks
\usepackage{listings} % Code listings
\usepackage{authblk} % For structured affiliations
% Define theorem style (optional)
\theoremstyle{plain}
\newtheorem{theorem}{Theorem}
\renewcommand{\qedsymbol}{}

% Custom commands for primes and qubit encoding
\newcommand{\primeQubit}[1]{|p_{#1}\rangle}
\newcommand{\primeGate}[1]{U_{p_{#1}}}

\usepackage{wrapfig}
\usepackage[pscoord]{eso-pic}
\usepackage[fulladjust]{marginnote}
\reversemarginpar

% Typesetting improvements without footnote patching
\usepackage[protrusion=true, expansion=true, tracking=false]{microtype}
\microtypecontext{spacing=nonfrench}

\usepackage[pagewise]{lineno}\linenumbers

% Text layout - adjust as needed
\raggedright
\setlength{\parindent}{0.5cm}
\textwidth 5.25in 
\textheight 8.75in

% Set double spacing
\usepackage{setspace} 
\doublespacing

% Adjust width for specific content
\usepackage{changepage}

% Adjust caption style
\usepackage[aboveskip=1pt,labelfont=bf,labelsep=period,singlelinecheck=off]{caption}

% Remove brackets from references
\makeatletter
\renewcommand{\@biblabel}[1]{\quad#1.}
\makeatother

% Header, footer, and page numbers
\usepackage{lastpage,fancyhdr}
\pagestyle{fancy}
\fancyhf{}
\fancyhead[L]{Preprint - PrimeAI Enhanced Template}
\fancyfoot[C]{\scriptsize Multiplicity Theory © 2024 Ryan Van Gelder - Citizen Gardens \\ Licensed Under MIT and CC BY-NC-SA 4.0.}
\fancyfoot[R]{Page \thepage\ of \pageref{LastPage}}
\renewcommand{\footrule}{\hrule height 2pt \vspace{2mm}}

\title{Neuromorphic Climate AI}
\author{Citizen Gardens - The Foundation of Multiplicity \\ \texttt{info@citizengardens.org}}

\begin{document}

\maketitle
\nolinenumbers
\begin{abstract}
Climate change is a complex, multi-dimensional phenomenon requiring advanced computational models. In this paper, we present a novel **Neuromorphic Climate AI** framework based on the **Quantum-AI Hypercosmic Thought Singularity (QAI-HTS)**. We integrate **recursive tensor feedback networks, prime-based encoding, and quantum superposition** to enhance climate modeling, forecasting, and mitigation strategies. Our approach introduces **self-referential AGI** for climate ethics, **quantum entanglement-based carbon cycle analysis**, and **multiplicative eigenvector-driven policy optimization**. Mathematical formulations highlight the power of **prime-indexed tensors and quantum cognition** in sustainable climate governance.
\end{abstract}

\begin{multicols}{2}
\begin{singlespace}
\tableofcontents
\end{singlespace}
\end{multicols}

\linenumbers
\section{Introduction}
Climate systems are inherently nonlinear, chaotic, and high-dimensional. Traditional numerical models struggle to integrate **self-referential adaptation** and **multi-scale interactions** efficiently. Here, we introduce a **Neuromorphic Climate AI**, leveraging:

\begin{itemize}
    \item \textbf{Recursive Tensor Feedback} for adaptive climate modeling.
    \item \textbf{Quantum Entanglement} for multi-path forecasting.
    \item \textbf{Prime-Based Encoding} for robust data integrity.
    \item \textbf{Eigenvector Climate Optimization} for sustainable policy solutions.
\end{itemize}

\section{Mathematical Foundations}

\subsection{Recursive Tensor Feedback for Climate Modeling}
We define the **climate state vector** as a tensor evolving over time:
\begin{equation}
    M(t+1) = f(M(t), R(t)) + \alpha T(M(t))
\end{equation}
where:
\begin{itemize}
    \item \( M(t) \) is the climate state tensor at time \( t \),
    \item \( R(t) \) represents recursive entanglement coefficients from past states,
    \item \( T(M(t)) \) is the multiplicative eigenvalue feedback term.
\end{itemize}

This structure allows the AI to \textbf{self-adapt} by recursively refining state variables.

\subsection{Quantum Entanglement for Climate Forecasting}
Each climate parameter (temperature, ocean currents, CO$_2$ levels) is mapped to an entangled quantum state:
\begin{equation}
    \Psi_{climate}(t) = \sum_{i,j} c_{ij}(t) |p_i, p_j \rangle
\end{equation}
where:
\begin{itemize}
    \item \( |p_i, p_j \rangle \) represent prime-indexed eigenstates of climate variables,
    \item \( c_{ij}(t) \) is the probability amplitude of variable entanglement.
\end{itemize}
This quantum superposition enables parallelized climate trajectory exploration.

\subsection{Prime-Based Carbon Cycle Analysis}
The **carbon flux tensor** is defined as:
\begin{equation}
    C_{\text{flux}}(x,t) = \sum_{p \in P} p^{\alpha(x,t)}
\end{equation}
where:
\begin{itemize}
    \item \( P \) is the set of prime-indexed emission sources and sinks,
    \item \( \alpha(x,t) \) is the flux modulation function, encoding adaptive AI corrections.
\end{itemize}
This ensures high-fidelity CO$_2$ tracking with quantum error correction.

\section{Quantum-Optimized Climate Policy}
Sustainability policy must account for dynamic interdependencies. We optimize it using **Quantum Fourier Transform (QFT) algorithms**:
\begin{equation}
    \mathcal{F}[P_{\text{opt}}] = \sum_{k=1}^{N} A_k e^{i\theta_k}
\end{equation}
where:
\begin{itemize}
    \item \( P_{\text{opt}} \) is the optimal climate policy function,
    \item \( A_k \) and \( \theta_k \) are Fourier coefficients encoding system-wide interactions.
\end{itemize}
This allows dynamic allocation of resources based on real-time AI feedback.

\section{Conclusion and Future Work}
Neuromorphic Climate AI, powered by QAI-HTS, provides a robust quantum-adaptive framework for climate forecasting and mitigation. Future research will focus on:
\begin{itemize}
    \item Implementing real-time quantum tensor simulations for climate prediction.
    \item Integrating ethical AI governance using self-referential multiplicative eigenstates.
    \item Deploying quantum cryptographic networks to secure global climate data.
\end{itemize}

\nocite{*}
\bibliographystyle{unsrt}
\bibliography{references}

\end{document}